\section{準備}\label{sec:prelim}

\subsection{変数}
KN~\cite{Kameyama17} に従い，$q=e^{i\epsilon_1}$，$t=e^{i\epsilon_2}$ を $\Omega$ 背景 $(z_1,z_2)\mapsto(qz_1,t^{-1}z_2)$ の変数とし，
$SU(N)$ の不変量の large $N$ 極限を $a=t^N$ で表す．
太字の変数 $(\ba,\bq,\bt)$ を
\begin{equation}\label{eq:cov}
a=-\ba^2\bt,\qquad q^{1/2}=-\bq\bt,\qquad t^{1/2}=\bq
\end{equation}
で定める（KN の式 (cov)）．これは Dunfield--Gukov--Rasmussen~\cite{Dunfield05} の superpolynomial の変数と同じであり，
$t=q$ は $\bt=-1$ に当たる．Aganagic--Shakirov~\cite{Aganagic11} の太字の変数も同じである．
Cherednik--Danilenko~\cite{Cherednik14} の DAHA の変数 $(q,t,a)$ は KN の $(q,t)$ と同じで，$a_{\mathrm{DAHA}}=-a$ である．
半整数冪の分枝は \eqref{eq:cov} で固定する．$\hf_\rho$ の全体の符号は規約であり，最低次の項が正になる向きにとる．

\subsection{Macdonald 多項式と自明な結び目}
$P_\lambda(x;q,t)$ を monic な Macdonald 多項式とし，内積 $\langle p_\lambda,p_\mu\rangle_{q,t}=\delta_{\lambda\mu}z_\lambda\prod_i(1-q^{\lambda_i})/(1-t^{\lambda_i})$ に関する
$g_\lambda=(q/t)^{|\lambda|}/\langle P_\lambda,P_\lambda\rangle_{q,t}$ を用いる．
自明な結び目の不変量は主特殊化
\[
P_\lambda(1,t,\dots,t^{N-1})=\prod_{s\in\lambda}\frac{t^{l'(s)}-a\,q^{a'(s)}}{1-q^{a(s)}t^{l(s)+1}}\qquad(a=t^N)
\]
を対称化したものである．$a(s),l(s)$ は腕と脚の長さ，$a'(s),l'(s)$ は余腕と余脚の長さである．
$n(\lambda)=\sum_i(i-1)\lambda_i$ とおき，
\[
T_\lambda=q^{n(\lambda')}\,t^{-n(\lambda)}
\]
を refined の framing の固有値と呼ぶ．$t=q$ では $q^{\sum_{s\in\lambda}c(s)}$（$c(s)$ は content）になる．

\subsection{refined CS 不変量}
KN と同じく，unreduced の refined CS 不変量を $\overline{\rCS}_\lambda(K;a,q,t)$，reduced の不変量を $\rCS_\lambda(K;a,q,t)$ と書く．両者は
\[
\overline{\rCS}_\lambda(K;a,q,t)=\overline{\rCS}_\lambda(\bigcirc;a,q,t)\,\rCS_\lambda(K;a,q,t)
\]
で結ばれる．ここで $\bigcirc$ は自明な結び目である．
トーラス結び目では，Aganagic--Shakirov の $S,T$ 行列による定義~\cite{Aganagic11}，Cherednik の DAHA による定義~\cite{Cherednik11}，
Gorsky--Neguţ~\cite{Gorsky13b} の楕円 Hall 代数による表示が知られている．
本論文では後者で計算する（\S\ref{sec:inputs}）．
これらの定義は単項式倍を除いてしか値を決めない．KN の BPS 抽出は単項式倍で結果が変わるので，単項式の正規化を \S\ref{sec:normalization} で固定する．
$\U(N)$ と $\SU(N)$ の違いは framing と linking number の補正だけで，本論文では区別しない（KN の脚注と同じ扱い）．

\subsection{refined reformulated invariants}
KN は，$q$-brane と $\bar t$-brane の二つの配位について，幾何学的遷移の等式
\begin{align}
\sum_\lambda \overline{\rCS}_\lambda(K;a,q,t)\,g_\lambda(q,t)\,P_\lambda(x;q,t)
&=\exp\Bigl(\sum_{d=1}^\infty\sum_\mu \frac1d\,\frac{f^q_\mu(K;a^d,q^d,t^d)}{q^{\frac d2}-q^{-\frac d2}}\,s_\mu(x^d)\Bigr),\label{eq:GT1}\\
\sum_\lambda \overline{\rCS}_\lambda(K;a,q,t)\,P_{\lambda^T}(-x;t,q)
&=\exp\Bigl(\sum_{d=1}^\infty\sum_\mu \frac1d\,\frac{f^{\bar t}_\mu(K;a^d,q^d,t^d)}{t^{-\frac d2}-t^{\frac d2}}\,s_\mu(x^d)\Bigr)\label{eq:GT2}
\end{align}
を提案した．
ここから定まる refined reformulated invariants $f^q_\mu$, $f^{\bar t}_\mu$ は，行列 $M_{\mu\rho}(t)$ により
\begin{equation}\label{eq:f-to-hatf}
f^q_\mu(K;a,q,t)=\sum_\rho M_{\mu\rho}(t)\,\wh f_\rho(K;a,q,t),\qquad f^{\bar t}_\mu(K;a,q,t)=\sum_\rho M_{\mu\rho}(q^{-1})\,\wh f_\rho(K;a,q,t)
\end{equation}
と書け，二つの式が同じ $\wh f_\rho$ を与えることが KN の枠組みの整合性になる．
$f^q_\mu$ の尺度は KN の付録 B の式に合わせる．本文の \eqref{eq:GT1} をそのまま展開した値とは $(q/t)^{|\mu|/2}$ だけ違う（付録~\ref{app:typos}）．

KN の予想は，$a$ 次数を $\pm\frac12$ ずらした後で $\wh f_\rho$ が
\begin{equation}\label{eq:hat-f-2}
\wh f_\rho(K;a,q,t)=\sum_{g,\beta,J_r}(-1)^{2J_r}\,\wh N_{\rho,g,\beta,J_r}(K)\,(q^{\frac12}-q^{-\frac12})^{g}(t^{-\frac12}-t^{\frac12})^{g}\Bigl(\frac qt\Bigr)^{J_r-\frac\beta2}a^{\beta}
\end{equation}
の形をとり，$\wh N_{\rho,g,\beta,J_r}(K)$ が非負整数になる，というものである．
$a$ 次数のずらしは，$|\rho|=1$ で $a^{-1/2}$，$|\rho|=3$ で $a^{1/2}$ とする．
\eqref{eq:cov} の変数では，\eqref{eq:hat-f-2} は
\[
\wh f_\rho=\sum_{g,\beta,J_r}\wh N_{\rho,g,\beta,J_r}\,(\bq\bt-\bq^{-1}\bt^{-1})^{g}(\bq-\bq^{-1})^{g}\,\ba^{2\beta}\,\bt^{2J_r}
\]
となる（KN の $\wh{\mathbf N}_{\rho,g,\beta,F}$ で $F=2J_r$ としたもの）．
$G=(\bq\bt-\bq^{-1}\bt^{-1})(\bq-\bq^{-1})$ とおき，この展開を $G$ 展開と呼ぶ．
$G$ 展開が整数係数で閉じるのは，$\ba^{2\beta}$ の各行が
\begin{equation}\label{eq:sigma}
\sigma:\ (e_\ba,e_\bq,e_\bt)\mapsto(e_\ba,-e_\bq,e_\bt-e_\bq)
\end{equation}
で不変で，$\bq$ の指数が偶数のときである．
KN は，展開をする前の素朴な変数変換の係数
\[
\wh f_\rho(K;a=-\ba^2\bt,\,q=\bq^2\bt^2,\,t=\bq^2)=\sum_{i,j,k}\wh{\mathbf N}_{\rho;i,j,k}(K)\,\ba^{2i}\bq^{2j}\bt^{k}
\]
も非負整数になることを観察した．
本論文で refined LMOV 不変量と呼ぶのは，$G$ 展開の前の $\wh{\mathbf N}_{\rho;i,j,k}$ と，後の $\wh N_{\rho,g,\beta,J_r}$ の二つである．$f^q$, $f^{\bar t}$ の係数は付録で補足として扱う．

\begin{definition}[正値性]\label{def:positivity}
$\wh f_\rho$ が\emph{正}であるとは，全体の符号を一つ選ぶと，$\wh{\mathbf N}_{\rho;i,j,k}$ と $\wh N_{\rho,g,\beta,J_r}$ の全てが非負整数になることをいう．
\end{definition}

$\yng(1)$ では $\hf_{\yng(1)}=(a^{1/2}-a^{-1/2})\rCS_{\yng(1)}$ である．
したがって，$\rCS_{\yng(1)}$ を太字の変数で書いた $P$ に対し，$\hf_{\yng(1)}$ は（a 次数をずらした後）$(1+\ba^2\bt)\ba^{-2}\bt^{-1}P$ になる．
$P$ が HOMFLY-PT ホモロジーの Poincaré 多項式として解釈されるなら，$\yng(1)$ の正値性は自然に期待される．
